Everything is read in Gauge Theory Gravity on one flat chart: the cut's induced metric, its Gauss curvature and its rotation-gauge field are read from the gravitational field $(\hb, \Omega)$ through Fact \claimref{synthesis.imported-gtg-dictionary}; $\omega^{ab}$ and $R^{ab}$ are $\Omega$ and its field strength written as forms in an orthonormal physical frame, as in Proposition \claimref{dark-energy.wall-counterterm}. By Theorem \claimref{synthesis.euler-term-horizon-entropy-shift}(i) and Fact \claimref{synthesis.imported-wald-noether-charge-entropy}(c) the Jacobson-Myers density is $\tfrac14(1 + 2\alpha\,{}^{(2)}\!R) = \tfrac14 + \alpha K$, so $S[P] = A(P)/4 + \alpha\int_{P}K\,dA$; this is (W$_{\mathrm{loc}}$) of Theorem \claimref{synthesis.series-wald-entropy-changes-without-flux-on-local-horizons}.

\emph{(i)} In two dimensions the rotation group is abelian and the structure equations of an orthonormal coframe $(e^{2}, e^{3})$ read $de^{2} = -\omega_{\Sigma}\wedge e^{3}$, $de^{3} = \omega_{\Sigma}\wedge e^{2}$, with curvature $d\omega_{\Sigma} = K\,e^{2}\wedge e^{3}$ (the Gauss curvature; Cartan's structure equations, Fact \claimref{synthesis.imported-gauss-bonnet-with-corners-poincare-hopf-and-edge-term}(a)). This is the two-dimensional transgression: the Euler form of the cut is $d$ of its Chern-Simons form, which in two dimensions is the connection itself. Stokes' theorem on a patch $P$ on which the frame is smooth gives $\int_{P}K\,dA = \oint_{\partial P}\omega_{\Sigma}$, and $\oint_{\partial P}\omega_{\Sigma}$ is the holonomy angle, measured against the frame, of a vector carried around the edge by the cut's own Levi-Civita connection (its rotation-gauge field); this is the cut's intrinsic parallel transport, not the series' tractor transport. A rotation $e' = R(\theta)e$ gives $\omega'_{\Sigma} = \omega_{\Sigma} - d\theta$ (the so(2) conjugation is trivial and $R\,dR^{-1}$ supplies $-d\theta$), so the edge integral changes by $-\oint d\theta$: zero when $\theta$ is a function on $P$, $-2\pi w$ when $\theta$ winds $w$ times along the edge. The check takes cuts $ds^{2} + \rho(s)^{2}d\phi^{2}$ with five polynomial profiles $\rho$, computes $K$ from the metric through its Christoffel symbols and the Riemann tensor (rational functions of $s$), reads $\omega_{\Sigma} = -\rho'\,d\phi$ off the structure equations for the frame $(\partial_{s}, \rho^{-1}\partial_{\phi})$, and confirms $d\omega_{\Sigma} = K\rho\,ds\wedge d\phi$ identically, the equality of area and edge integrals on rectangles $[s_{1}, s_{2}]\times[0, \phi_{0}]$, and the frame rotation law with $\cos\theta = (1 - t^{2})/(1 + t^{2})$, $\sin\theta = 2t/(1 + t^{2})$ for three polynomials $t(s)$, where $d\theta = 2t'\,ds/(1 + t^{2})$.

For the link with the wall counterterm: for $Q = \epsilon_{abcd}(\omega^{ab}\wedge d\omega^{cd} + \tfrac23\omega^{ab}\wedge\omega^{c}{}_{e}\wedge\omega^{ed})$ and $\delta\omega = D\lambda$, the first variation is $\delta Q = 2\epsilon(\delta\omega, R) - d\,\epsilon(\omega, \delta\omega)$, and with $\epsilon(D\lambda, R) = d\,\epsilon(\lambda, R)$ (Bianchi identity and invariance of $\epsilon$) this is $d\,\epsilon(\lambda, d\omega)$, the standard gauge variation of a Chern-Simons form (Fact \claimref{synthesis.imported-chern-weil-and-chern-simons-forms}(c)); the check confirms $\delta_{\lambda}Q = d\,\epsilon(\lambda, d\omega)$ exactly as an identity of forms for a polynomial $\mathfrak{so}(1,3)$ connection with legs in all four directions and a polynomial parameter, with $\epsilon_{abcd}$ read from $e_{a}e_{b}e_{c}e_{d} = \epsilon_{abcd}I$ in $\Cl(1,3)$. For $\lambda = e_{0}\wedge e_{1}$, $\epsilon_{abcd}\lambda^{ab}d\omega^{cd} = 2\epsilon_{01cd}d\omega^{cd} = 4\,d\omega^{23}$. On the cut $x_{0} = x_{1} = 0$ in a frame with $e_{2}, e_{3}$ tangent, $R^{23} - d\omega^{23} = \omega^{2}{}_{0}\wedge\omega^{03} + \omega^{2}{}_{1}\wedge\omega^{13} = -\omega^{20}\wedge\omega^{30} + \omega^{21}\wedge\omega^{31}$, and substituting $\omega^{A0} = b_{A} + a_{A}/2$, $\omega^{A1} = b_{A} - a_{A}/2$ gives $-(b_{2}\wedge a_{3} + a_{2}\wedge b_{3})$; with $\eta_{11} = -1$, $a_{A}$ is (up to sign) the second fundamental form along $k = e_{0} + e_{1}$, so on an admissible cut $R^{23} = d\omega^{23}$ there (the Gauss equation; the check confirms the identity on a generic connection and that the defect is nonzero when $a \neq 0$). So, as a reading supported by these identities, in an adapted frame on admissible cuts the frame-relative edge term of (i) is the wall counterterm's gauge variation under the cut's binormal, as the wall counterterm is the primitive of the Euler term on a wall: in both cases a transgression of a topological density, defined relative to a frame, with a topological residual.

\emph{(ii)} Gauss-Bonnet for a compact surface with piecewise smooth boundary (Fact \claimref{synthesis.imported-gauss-bonnet-with-corners-poincare-hopf-and-edge-term}(a)): $\int_{P}K\,dA + \oint_{\partial P}k_{g}\,ds + \sum_{j}\varepsilon_{j} = 2\pi\chi(P)$, with $k_{g}$ taken with respect to the inward normal and $\varepsilon_{j}$ the turning angles at corners. Multiplying by $\alpha$ gives the displayed identity, and $2\pi\alpha = -3\pi/(2\Lambda)$. The check computes $k_{g}$ of the circles $s = \mathrm{const}$ as $g(\nabla_{T}T, N)$ from the Christoffel symbols, $k_{g} = \mp\rho'/\rho$ for $N = \pm\partial_{s}$, verifies that the meridians are geodesics, and confirms the identity in units of $2\pi$ on rectangles (four right corners, total turn one), polar caps (no corners) and annuli ($\chi = 0$) for every profile. A local horizon moves its cuts along generators by a smooth deformation, so a patch keeps its topology and $\delta(2\pi\alpha\chi(P)) = 0$; hence $\delta S' = \delta A/4$. On the radial generators of a Schwarzschild-de Sitter hole, by Theorem \claimref{synthesis.series-wald-entropy-changes-without-flux-on-local-horizons}(ii), $K = 2M/r^{3} + \Lambda/3$ on the admissible cuts, $\delta A = 0$ and $\delta Q = 0$, while the bare density is $-3M/(2\Lambda r^{3})$; the check confirms $K = (1 - f)/r^{2}$, the bare density, the nonzero two-signed bare change, and that $\tfrac14 + \alpha K - \alpha K = \tfrac14$ at five points along each generator for three holes, both signs and two patch lengths. In the thin-pencil limit (H2) of Theorem \claimref{synthesis.einstein-equations-from-c-loc-and-clausius} the constant $2\pi\alpha\chi(P)$ is the same at both ends of the pencil, so it drops out of $\delta S'$ at every $a_{0}$; the first-order rate of Theorem \claimref{synthesis.euler-entropy-first-order-rate-on-local-patches}(i) is the $\lambda$-derivative of $\alpha\int K\,dA$, which by Gauss-Bonnet is minus the derivative of the edge term, so it is absent from $S'$ and the defect of (ii) of that theorem vanishes for every admissible element.

\emph{(iii)} Suppose $Y$ assigns to each oriented edge a number, with $Y(-c) = -Y(c)$, determined by the field and the cut, and $S_{E}[P] = -Y(\partial P)$ for every disk patch. Cut a closed cut $\Sigma$ of genus zero along a circle $c$ into disks $P_{1}, P_{2}$; then $S_{E}[\Sigma] = -Y(c) - Y(-c) = 0$, while $S_{E}[\Sigma] = \alpha\int_{\Sigma}K\,dA = 4\pi\alpha \neq 0$ (the check confirms $\int K\,dA = 4\pi$ on three closed profiles $\rho = s(L - s)(1 + c\,s(L - s))/L$, $\rho'(0) = 1$, $\rho'(L) = -1$, positive inside). So no such $Y$ exists. The geodesic-curvature term is additive on smooth edges but its corner turns are not (two rectangles glued along a full edge have eight right corners, their union four), which is where $2\pi\chi(P)$ sits. The frame term $\oint\omega_{\Sigma}$ is additive for one frame, but no tangent frame is smooth on a sphere (Poincar\'e-Hopf, Fact \claimref{synthesis.imported-gauss-bonnet-with-corners-poincare-hopf-and-edge-term}(b)): on the caps $s < s_{e}$ and $s > s_{e}$ the frames rotated by $\mp\phi$ from the polar frame are smooth, with $\omega_{N} = (1 - \rho')\,d\phi$ and $\omega_{S} = (-1 - \rho')\,d\phi$ vanishing at their poles, and along every separating circle $\omega_{N} - \omega_{S} = 2\,d\phi$, winding $2 = \chi(S^{2})$ (the clutching number, Fact \claimref{synthesis.imported-gauss-bonnet-with-corners-poincare-hopf-and-edge-term}(b)); the check confirms both cap integrals as edge integrals and their sum $4\pi$ at three separating circles of three closed profiles. Hence on a closed cut every choice leaves $2\pi\alpha\chi = -3\pi\chi/(2\Lambda)$, which is the shift of Theorem \claimref{synthesis.euler-term-horizon-entropy-shift}(ii) ($-3\pi/\Lambda$ on spheres) and, with $3\pi/\Lambda$ the de Sitter entropy, $-\chi/2$ times it; Theorem \claimref{synthesis.sds-clausius-fixes-area-constant}(v) excludes it under (Z) or nonnegativity. On a disk patch, $S'[P] = A(P)/4 + 2\pi\alpha = A(P)/4 - 3\pi/(2\Lambda)$, so for $\Lambda > 0$, $S'[P] < 0$ exactly when $A(P) < 6\pi/\Lambda$, half the de Sitter horizon area $4\pi\cdot 3/\Lambda$ (the check confirms the threshold and the sign at three values of $\Lambda$); for $\Lambda < 0$ the constant is positive. Two disks glued along an edge have $S'$ summing to $A/4 - 3\pi/\Lambda$ against $A/4 - 3\pi/(2\Lambda)$ for their union, so $S'$ is not additive.

\emph{(iv)} By (ii), with the edge term, entropy changes of open patches are $\delta A/4$ for both readings, so (H1) holds for $S'$ exactly where it holds for $A/4$, and the bare functional's open-patch defect is the variation of the omitted edge term. By (iii) the remaining difference, $2\pi\alpha\chi(P)$ on every patch and $2\pi\alpha\chi$ on closed cuts, is one topological constant that no additive edge term removes. Under the integrated reading of (Loc) in Premise \claimref{synthesis.c-loc-local-horizon-premise} the entropy of a patch is $\int\eta\,dA$, which $A(P)/4 + 2\pi\alpha\chi(P)$ is not for $\alpha \neq 0$ (the constant does not scale with the patch), so that reading excludes the constant; on the differences reading the constant does not enter. Nothing in (i) to (iii) uses a property of the frozen point, the canonical tractor or the tractor gauge field beyond the value of $\alpha$, so transport does not fix the constant: the exclusion of the Euler term from horizon entropy remains the adopted decision. Whether the edge term of (ii) is the one the action prescribes on open patches is argued, not proved: it is the Gibbons-Hawking term of the two-dimensional Euler functional (Fact \claimref{synthesis.imported-gauss-bonnet-with-corners-poincare-hopf-and-edge-term}(c)), the term that makes its variation with the edge metric held fixed well posed, and Wald's Noether charge is fixed only up to closed forms, which on an open patch are edge terms; the Iyer-Wald and Jacobson-Kang-Myers rules (cited from memory) fix that freedom only on closed cuts.


\emph{Remark (argued, not used above).} In the BF variables of Theorem \claimref{synthesis.frozen-point-action-is-bf-transport-broken-by-the-blade-i-wedge-x}, the on-shell two-form field is proportional to the pseudoscalar dual of the Lorentz block of the Cartan curvature, and the pseudoscalar maps the cut's tangent bivector $e_{2}\wedge e_{3}$ to the binormal $e_{0}\wedge e_{1}$ up to sign; so the Jacobson-Myers density $-\tfrac{3}{4\Lambda}F(e_{2}, e_{3}, e_{2}, e_{3})$ of Theorem \claimref{synthesis.series-wald-entropy-changes-without-flux-on-local-horizons}(i) is, up to a constant whose sign is not fixed here, the flux of that field through the cut (the usual reading of the Wald entropy of a BF-type action, cited from memory). Splitting the Cartan curvature into its Lorentz part and its soldering part $-\tfrac{\Lambda}{3}e\wedge e$, the soldering part gives $-\tfrac{3}{4\Lambda}\cdot(-\tfrac{\Lambda}{3}) = \tfrac14$, the area density, and the Lorentz part gives the Euler part treated in (i) to (iii). The adopted reading of horizon entropy as (Loc) keeps the flux of the soldering part only.
